\section{Introduction}\label{sec:intro}

Temporal graphs model connections whose availability changes with time~\cite{holme2012,michail2016}. In arc routing, edges are tasks, while vehicle scheduling assembles time-ordered trips into duties~\cite{edmonds1973,corberan2015,bunte2009}. These settings motivate covering edges by walks that respect their time labels. We consider undirected edges, choosing their traversal direction together with the walks. Experiments ignore recorded directions and travel times, testing correctness and scalability rather than a routing application.

An exact edge cover partitions temporal edges among walks, traversing each edge once. The input prescribes start counts and, in \FIX{}, end counts without pairing individual starts and ends. In \FREE{}, the solution chooses the end vertices. Empty walks are allowed. A non-strict walk uses non-decreasing time labels and may traverse several edges of one snapshot, whereas a strict walk uses strictly increasing labels and at most one edge per snapshot. Equal time stamps occur in coarsely recorded contact data, and the strict/non-strict distinction can change complexity~\cite{casteigts2024,erlebach2023}, including for temporal exploration.

Deligkas et al.~\cite{deligkas2025} introduced temporal exact edge cover at AAAI 2025, establishing the problem and its complexity results. This paper builds on their problem and notation. They claimed polynomial solvability for non-strict walks through a static reduction. Their NP-completeness result without prescribed starts concerns minimizing walk count. Strict walks with prescribed terminals were left open. The proof sketch of their Theorem~6.3 misses a temporal availability condition: every snapshot component containing an edge must contain a walk immediately before its time, which static connectivity to a terminal does not ensure. Section~\ref{sec:counterexample} gives a five-vertex instance whose static graph has an exact edge cover but whose temporal graph does not, under the end-terminal construction in Figure~3 of~\cite{deligkas2025}. This gap does not refute their claim. We prove the non-strict polynomial result differently and settle the strict variants.

Our method characterizes realizable orientations by non-negative walk counts and a walk in every snapshot component. Only components with even degrees and counters can violate the latter. Pair counting exposes M-convex component states, and removing the dead state preserves M-convexity. This yields an M-convex intersection formulation and an augmenting path algorithm that repairs dead components or certifies infeasibility. To our knowledge, dead-state treatment is new. 

\paragraph{Contributions.}
\begin{enumerate}
\item \emph{Characterization and gap.} An orientation is realizable if and only if the counting condition holds and every snapshot component contains a walk (Theorem~\ref{thm:char}). Only even components can violate the latter (Lemma~\ref{lem:noneven}), and instances without them reduce to one maximum flow (Corollary~\ref{cor:flow}). Under the end-terminal construction shown in Figure~3 of~\cite{deligkas2025}, a five-vertex instance shows a gap in the static reduction (Proposition~\ref{prop:aaai}, Lemma~\ref{lem:aaai}).
\item \emph{M-convex feasibility algorithm.} We prove that removing the dead state preserves M-convexity for even-component states (Lemma~\ref{lem:gamma}(b)), yielding a polynomial algorithm via feasibility of an M-convex intersection (Theorem~\ref{thm:reduction}). It returns walks or an infeasibility certificate (Theorem~\ref{thm:correct}) and has the bound of Theorem~\ref{thm:time}.
\item \emph{Strict-walk hardness.} The strict problem is NP-complete with prescribed ends and six time labels, and with free ends and two time labels, empty walks allowed (Theorem~\ref{thm:strict}).
\end{enumerate}
